
\textbf{Lemme 1.5.2.} \uline{Soient} @@M0@@ \uline{un anneau}, @@M1@@ \uline{une} @@M2@@-\uline{algèbre de présentation finie}, @@M3@@ \uline{un} @@M4@@-\uline{module de présentation finie}. \uline{On suppose} @@M5@@ \uline{et} @@M6@@ \uline{plats sur} @@M7@@, \uline{et} \uline{que pour tout corps résiduel} @@M8@@ \uline{de} @@M9@@, @@M10@@ \uline{est de dimension projective} @@M11@@ \uline{sur} @@M12@@ (\uline{ce qui est le cas si} @@M13@@). \uline{Sous ces conditions}, @@M14@@ \uline{admet une résolution de longueur} @@M15@@ \uline{par des} @@M16@@-\uline{modules projectifs de type fini} :
@@M17@@
Ce lemme résulte immédiatement de EGA IV 12.3.2.

\textbf{Proposition 1.6.} \uline{Soient} @@M18@@ \uline{un schéma quasi-compact}, @@M19@@ \uline{un} @@M20@@\uline{-Module localement libre de rang} @@M21@@, @@M22@@ \uline{son algèbre symétrique}, \uline{avec sa graduation habituelle}, @@M23@@ \uline{un} @@M24@@\uline{-Module gradué de présentation finie}, \uline{plat sur} @@M25@@. \uline{Alors il existe une résolution graduée} :
@@M26@@
\uline{les} @@M27@@ \uline{étant des} @@M28@@\uline{-Modules gradués qui sont localement libres de type fini} ( \uline{en tant que} @@M29@@\uline{-Modules gradués}).

Montrons d'abord que pour tout @@M30@@-Module gradué @@M31@@, satisfaisant aux conditions de la proposition, il existe un homomorphisme gradué surjectif~: @@M32@@, où @@M33@@ est un @@M34@@-Module localement libre de type fini.

Comme @@M35@@ est de type fini, il existe un entier @@M36@@ tel que @@M37@@ soit engendré sur @@M38@@ par ses composantes de degré @@M39@@. On pose alors :
@@M40@@
l'homomorphisme de @@M41@@ dans @@M42@@ se définit grâce à (1.2.1), et il est clair qu'il est surjectif.

D'après 1.5.2, @@M43@@ admet une résolution (non graduée) par des @@M44@@-Modules localement libres de type fini sur chaque ouvert affine. Par suite, si 
